\chapter{Was ist NOMAD, und ab wann geht es ohne Internet?}

\section{Kurz erklärt}
NOMAD ist ein kleiner Server, der Wissen bereithält: Nachschlagewerke, Wikipedia, Anleitungen, Karten und auf Wunsch einen Sprachassistenten (KI). Sie bedienen ihn im Browser, wie eine Internetseite. Der Unterschied: Die Seite liegt auf \emph{Ihrem} Rechner. Ist das Wissen einmal geladen, brauchen Sie kein Internet mehr.

\bild[0.8]{img-boot2}{Die Startseite von NOMAD heißt „Kommandozentrale“. Von hier aus erreichen Sie alle Funktionen.}

\section{Der wichtigste Satz}
\merke{NOMAD ist erst offline-fähig, wenn Sie die Inhalte \textbf{vorher mit Internet geladen} haben. Ein frisch eingerichtetes NOMAD ohne geladene Inhalte ist eine leere Bibliothek.}

\section{Was geht wann?}
Die Tabelle zeigt, welche Stufe Sie erreicht haben und was dann ohne Internet möglich ist.

\renewcommand{\arraystretch}{1.35}
\begin{center}
\small
\begin{tabularx}{\linewidth}{>{\RaggedRight\arraybackslash}p{3.4cm}>{\RaggedRight\arraybackslash}X>{\RaggedRight\arraybackslash}X}
\toprule
\textbf{Stufe} & \textbf{Das geht ohne Internet} & \textbf{Das braucht Internet} \\
\midrule
\textbf{A. Gestartet} (USB-SSD gebootet, NOMAD eingerichtet) &
Rechner starten, Kommandozentrale öffnen, deutsche Anleitungen (Docs) lesen. Die Einrichtung selbst braucht kein Internet. &
Alles Weitere: Inhalte, Karten, KI-Modelle, Aktualisierungen. \\
\addlinespace
\textbf{B. Inhalte geladen} (Kapitel~5) &
Wikipedia, Nachschlagewerke, Anleitungen durchsuchen und lesen (Kiwix). Dies ist die Stufe, die Sie für den Notfall wollen. &
Neue Inhalte, neue Fassungen, Updates von NOMAD. \\
\addlinespace
\textbf{C. Karten geladen} &
Kartenansicht der geladenen Länder (zum Beispiel Deutschland, ab \SI{0,1}{GB}). &
Weitere Länder laden. \\
\addlinespace
\textbf{D. KI eingerichtet} (Kapitel~10) &
Fragen an den lokalen Sprachassistenten, ohne dass etwas ins Internet geht. &
Das Sprachmodell einmal herunterladen. \\
\bottomrule
\end{tabularx}
\end{center}

\hinweis{Eselsbrücke: Alles, was Sie \textbf{herunterladen}, braucht Internet. Alles, was Sie \textbf{lesen oder suchen}, nicht.}

\section{Wie viel Platz brauchen die Inhalte?}
NOMAD selbst belegt etwa 5~GB. Der Rest hängt davon ab, was Sie laden. Die Zahlen stammen aus den Katalogen von NOMAD (Stand Oktober 2026); die Stufen \emph{Standard} und \emph{Umfassend} enthalten jeweils die darunter liegende Stufe mit. NOMAD zeigt Größen mit Faktor 1024 an, deshalb stehen dort Zahlen, die rund 2\,\% kleiner sind als hier.

\textbf{Deutsche Inhalte (empfohlen):}

\renewcommand{\arraystretch}{1.25}
\begin{center}
\small
\begin{tabular}{@{}lrrr@{}}
\toprule
\textbf{Themenbereich} & \textbf{Basis} & \textbf{Standard} & \textbf{Umfassend} \\
\midrule
Medizin (Deutsch) & \SI{0,16}{GB} & \SI{0,20}{GB} & \SI{0,69}{GB} \\
Nachschlagen und Lernen (Deutsch) & \SI{0,26}{GB} & \SI{6,3}{GB} & \SI{18,6}{GB} \\
Alltag und Technik (Deutsch) & \SI{0,03}{GB} & \SI{0,21}{GB} & \SI{3,8}{GB} \\
\midrule
\textbf{Summe der deutschen Bereiche} & \textbf{\SI{0,45}{GB}} & \textbf{\SI{6,7}{GB}} & \textbf{\SI{23,1}{GB}} \\
\bottomrule
\end{tabular}
\end{center}

Dazu kommt die \textbf{deutsche Wikipedia} nach Wunsch: von \SI{137}{MB} (die wichtigsten Artikel, kurz) über \SI{1,2}{GB} (Beliebte Artikel) und \SI{18}{GB} (alles ohne Bilder) bis \SI{51}{GB} (alles mit Bildern); mehr dazu in Kapitel~5.

\textbf{Englische Inhalte} gibt es zusätzlich, vor allem für Überleben, Vorsorge und Medizin-Nachschlagewerke (Kategorien mit dem Zusatz „Englisch“). Alle englischen Bereiche zusammen sind \SI{6,3}{GB} (Essential), \SI{35,7}{GB} (Standard) und \SI{89,1}{GB} (Comprehensive).

\section{Wie lange dauert das Laden?}
Die Dauer hängt von Ihrem Internetanschluss ab. Die Tabelle zeigt die reine Übertragungszeit bei voller Geschwindigkeit:

\begin{center}\small
\begin{tabular}{@{}lrrr@{}}
\toprule
\textbf{Menge} & \textbf{16 Mbit/s (DSL klein)} & \textbf{50 Mbit/s} & \textbf{250 Mbit/s} \\
\midrule
\SI{0,45}{GB} (Deutsch, Basis) & 4 min & 1 min & 15 s \\
\SI{6,7}{GB} (Deutsch, Standard) & 56 min & 18 min & 4 min \\
\SI{23}{GB} (Deutsch, Umfassend) & 3,2 Std. & 1,0 Std. & 12 min \\
\SI{89}{GB} (Englisch, alles) & 12,4 Std. & 4,0 Std. & 48 min \\
\bottomrule
\end{tabular}
\end{center}
In der Praxis dauert es eher 20 bis 50\,\% länger. Der Rechner darf währenddessen nicht in den Ruhezustand gehen (im NOMAD-System ist er abgeschaltet).

\section{Zwei Wege zum Ziel}
\begin{description}[leftmargin=1.2em,style=nextline]
\item[Weg 2 – USB-SSD (empfohlen)] Sie schreiben ein fertiges System auf eine USB-SSD und starten Ihren Windows-PC davon. Windows und Ihre Daten auf der Festplatte des PCs bleiben unberührt. Siehe Kapitel~2.
\item[Weg 1 – eigener Rechner mit Ubuntu] Sie installieren Ubuntu Linux auf einem Rechner, der danach nur noch NOMAD betreibt (zum Beispiel ein altes Notebook). Siehe Kapitel~3 und~4.
\end{description}
